\section{Six positive-helicity photons}\label{sec:eight-point}
For $\M(1_e,2_E,3^+,4^+,5^+,6^+,7^+,8^+)$ there are $6!=720$
orders of the photons along the fermion line. An electron propagator
separates a nonempty proper subset of the photons from its complement,
giving $2^6-2=62$ distinct physical channels. The leg containing the
electron distinguishes the two sides. Write $\sum_{S_6}$ for the
unnormalized sum over the six photon labels, acting on every occurrence
of each label, and set
\[
 C=m(\sqrt{2}e)^6\ang{\mathbf1\mathbf2},\qquad
 x_i=\frac{[i12i]}{\tau_{2i}} .
\]

\paragraph{Physical poles.}
Choose photons 3 and 4 as the first unordered pair. The first term
below is symmetric in all four remaining photons. The mass correction
selects a second pair, 5 and 6. Two contractions enter that correction,
\[
 U=[57][715\rangle,\qquad
 V=[8|p_1(p_5+p_6+p_7)|8].
\]
Splitting the three null momenta gives
$V=[58][815\rangle+[68][816\rangle+[78][817\rangle$.
With complementary labels shortened by momentum conservation, we find
\begin{equation}\label{eq:eight-physical}
\begin{aligned}
 \M_{++++++}^{\phys}={}&\frac{C}{48}\sum_{S_6}
 \frac{[34]^2}{\tau_{15}\tau_{16}\tau_{17}\tau_{18}}
 \bigg\{\frac{x_5 x_6 x_7 x_8}{\tau_{13}\tau_{14}}\\
 &+\frac{6m^2[56]^2}{\tau_{23}\tau_{24}\tau_{156}}
 \left[-x_7x_8+\frac{m^2[78]^2}{\tau_{234}}
 +\frac{4U}{\tau_{1567}}\left(x_8-\frac{V}{\tau_{234}}\right)\right]\bigg\}.
\end{aligned}
\end{equation}
The leading term has $\binom62=15$ distinct images.
It occurs $2!4!=48$ times under the full sum, accounting for its
coefficient. On distribution the entire expression has seven photon
permutation orbits.

\paragraph{Spurious poles.}
Let $H_i=[i13\rangle$ and $B=6m^2[78]^2/\tau_{234}$.
The eight orbits of the spurious-pole expression combine into
\begin{equation}\label{eq:eight-spurious}
\begin{aligned}
 \M_{++++++}^{\spur}={}&\frac{C}{24}\sum_{S_6}
 \frac{[34]}{\tau_{17}\tau_{18}\tau_{23}\ang{35}\ang{36}}
 \bigg\{\frac{H_5H_6}{\ang{34}\tau_{15}\tau_{16}}
 \left[\frac{H_7 H_8}{\ang{37}\ang{38}}-B\right]\\
 &+\frac{2B}{\tau_{1578}}\left([46]+\frac{[623\rangle}{\ang{34}}\right)
 \left[\frac{H_5}{\tau_{15}}+
       \frac{H_5-2\ang{37}[57]}{\tau_{178}}\right]\bigg\}.
\end{aligned}
\end{equation}

A different arrangement gives a single product for each photon order.
Take 7 as the reference for the four emissions $3,4,5,6$ and 8 as
its partner. The successive electron momenta enter the four numerators.
In the last emission, momentum conservation replaces
$p_1+p_3+p_4+p_5$ by $-(p_2+p_8)$ inside the sandwich, since $p_6$
and $p_7$ annihilate its endpoints. With
$\langle i(ab\ldots)j]=\langle i|(p_a+p_b+\cdots)|j]$, this gives
\begin{equation}\label{eq:eight-ordered}
 \M_{++++++}^{\spur}=-C\sum_{S_6}
 \frac{[78]\,\langle713]\,\langle7(13)4]\,
        \langle7(134)5]\,\langle7(28)6]}
 {\ang{78}\ang{73}\ang{74}\ang{75}\ang{76}\,
  \tau_{27}\tau_{13}\tau_{134}\tau_{1345}\tau_{278}} .
\end{equation}
Each choice of the ordered pair $(7,8)$ leaves $4!$ emission orders.
Thus $6\cdot5\cdot4!=720$ counts every term once. Equation~\eqref{eq:eight-ordered}
is the shortest to write. Equation~\eqref{eq:eight-spurious} has a lower
weighted complexity when the momentum sums are counted.

Both constructions obey the 62 signed factorization conditions and the
fixed-angle $E^{-4}$ bound. We also transcribed
Eqs.~\eqref{eq:eight-physical}--\eqref{eq:eight-ordered} and tested each
against the signed sum of 720 Feynman diagrams at 100,000 phase-space
points. Two points required 160-bit kinematics, where all three formulas
agreed within $9.1\times10^{-47}$ in relative amplitude. At fixed photon
helicity and summed over the two fermion spins, polarized MadGraph agreed
at all 100,000 test points. The largest relative differences in the squared
amplitude were $2.52\times10^{-12}$, $1.12\times10^{-11}$ and
$6.88\times10^{-12}$ for the three equations, respectively.
